%! Author = denis
%! Date = 20/08/2026

\begin{frame}{Circuitos Combinacionais x Sequenciais}

    \begin{columns}[T]

        \begin{column}{0.48\textwidth}

            \begin{block}{Combinacionais}
                A saída depende somente das entradas atuais.

                \[
                    Y=f(A,B,C)
                \]

                Exemplos:

                \begin{itemize}
                    \item Somadores
                    \item Comparadores
                    \item Multiplexadores
                    \item Decodificadores
                \end{itemize}
            \end{block}

        \end{column}

        \begin{column}{0.48\textwidth}

            \begin{block}{Sequenciais}
                A saída depende das entradas e do estado anterior.

                \[
                    Y=f(A,B,C,Q)
                \]

                Exemplos:

                \begin{itemize}
                    \item Flip-flops
                    \item Registradores
                    \item Contadores
                    \item Máquinas de estados
                \end{itemize}
            \end{block}

        \end{column}

    \end{columns}

\end{frame}